\section{Contexto y problema}

La carrera de Software de la Universidad Técnica Estatal de Quevedo realiza
pre-sustentaciones como parte de Desarrollo de Titulación II. En este proceso
participan el estudiante, el docente de la asignatura, el director del trabajo,
la coordinación de titulación y un tribunal formado por un presidente y dos
miembros. Durante el segundo periodo 2024--2025 se registraron 12 estudiantes;
dos se retiraron y diez llegaron a la predefensa. Por tanto, el volumen habitual
es de 10 a 12 estudiantes por periodo.

Actualmente, varias actividades se coordinan de forma manual y mediante correo
electrónico. Por ese medio se solicitan y entregan las actas, mientras que la
firma en papel retrasa el cierre del proceso y deja abierta la posibilidad de
extravíos. También se recordó un cambio repentino de tutor que produjo
conflictos porque no quedó registrado. El sistema propuesto debe conservar una
bitácora de esos cambios y mostrar el estado de cada solicitud a las personas
autorizadas.

Cada estudiante presenta dos o tres documentos PDF: el informe del trabajo, el
reporte de URKUND que incorpora el director y, cuando el tribunal formula
observaciones, una versión corregida. Con un máximo previsto de 20 MB por
archivo, el volumen de documentos entregados no supera los 720 MB para 12
estudiantes, sin contar las actas generadas. La carga se concentra en los días
anteriores a las predefensas y durante la semana 19 del décimo semestre.

La calificación y el acta son datos críticos. La nota final combina el 40\% del
tribunal y el 60\% del docente de Desarrollo de Titulación II, y luego se
registra en el sistema informático de la UTEQ. El porcentaje de URKUND también
afecta el resultado: por debajo del 10\% el trabajo se considera aceptable,
mientras que un valor superior al 20\% interrumpe la predefensa. Un error o una
versión desactualizada de estos datos puede cambiar la condición de aprobación
del estudiante.

Los principales interesados esperan resultados distintos: el estudiante
necesita entregar documentos y consultar el estado; coordinación requiere
organizar solicitudes, tutores, tribunales y reportes; el director debe adjuntar
el reporte de similitud; y el tribunal necesita revisar el expediente, registrar
su evaluación y completar el acta. Los avisos por correo, los listados, los
reportes y el estado visible pueden aceptar una actualización no inmediata,
pero las notas y las actas deben mantenerse íntegras y actualizadas.

La información de esta sección procede del resumen reconstruido por el equipo
después de la entrevista del 19 de noviembre de 2026. La Ing. Nancy Rodríguez
confirmó el contenido del resumen utilizado como fuente.
